%%%%%%%%%%%%%%%%%%%%%%%%%%%%%%%%%%%%%%%%%%%%%%%%%%%%%%%%%%%%%%%%%%%%%%%%%%%%%%%%%%%%%%%%%%%%%%%%%%%%%%%%%%%%%%%%%%%%%%%%
\section{Project closing}
\label{sec:closing}
%%%%%%%%%%%%%%%%%%%%%%%%%%%%%%%%%%%%%%%%%%%%%%%%%%%%%%%%%%%%%%%%%%%%%%%%%%%%%%%%%%%%%%%%%%%%%%%%%%%%%%%%%%%%%%%%%%%%%%%%
Project closing is the final phase of the project lifecycle. Its purpose is to formally conclude the collaboration, 
verify the completion of the project objectives and deliverables, document the project outcomes and ensure that all 
administrative outstanding tasks are completed. Project closing consists of two stages: the closing meeting and the end 
report.

%%%%%%%%%%%%%%%%%%%%%%%%%%%%%%%%%%%%%%%%%%%%%%%%%%%%%%%%%%%%%%%%%%%%%%%%%%%%%%%%%%%%%%%%%%%%%%%%%%%%%%%%%%%%%%%%%%%%%%%%
\subsection{Closing meeting}
\label{sec:closing:meeting}
%%%%%%%%%%%%%%%%%%%%%%%%%%%%%%%%%%%%%%%%%%%%%%%%%%%%%%%%%%%%%%%%%%%%%%%%%%%%%%%%%%%%%%%%%%%%%%%%%%%%%%%%%%%%%%%%%%%%%%%%
The closing meeting marks the conclusion of the collaboration between the eScience Center and the Lead Applicant. It is 
typically held as the last regular project meeting~\ref{sec:exec:team-meetings} during the last weeks of the project. 
The meeting provides an opportunity to review the project outcomes, discuss any remaining activities, agree on 
responsibilities after the project ends, and prepare the end report.

Typical discussion topics include:
\vspace{-0.2cm}
\begin{itemize}
 \item review of the project objectives and deliverables,
 \item outstanding publications, software releases or other project outputs,
 \item remaining knowledge transfer activities,
 \item software sustainability and maintenance plans,
 \item sharing the end report template, and
 \item opportunities for future collaboration.
\end{itemize}
\vspace{-0.2cm}

The conclusions and action points are recorded in the project log.

\textbf{Responsible: Lead RSE.} 

%%%%%%%%%%%%%%%%%%%%%%%%%%%%%%%%%%%%%%%%%%%%%%%%%%%%%%%%%%%%%%%%%%%%%%%%%%%%%%%%%%%%%%%%%%%%%%%%%%%%%%%%%%%%%%%%%%%%%%%%
\subsection{End report}
\label{sec:closing:end}
%%%%%%%%%%%%%%%%%%%%%%%%%%%%%%%%%%%%%%%%%%%%%%%%%%%%%%%%%%%%%%%%%%%%%%%%%%%%%%%%%%%%%%%%%%%%%%%%%%%%%%%%%%%%%%%%%%%%%%%%
The end report documents the outcomes of the project and provides the information required for the formal closure of 
the project. For internal projects, the end report is written by the LA using the template available in the 
\textit{Project templates} folder of the project portfolio and previously introduced by the Lead RSE at the closing 
meeting. 

\let\myhcolw\relax 
\newlength{\myhcolw}
\setlength{\myhcolw}{0.8\textwidth}
\begin{table}[!h]
 \captionsetup{
   labelformat=fancytable,
   labelsep=space,
   labelfont=bf,
   textfont=bf,
   justification=centering,
   singlelinecheck=true
 }
 \caption{End report process overview}
 \begin{booktabs}{
    colspec={|>{\bfseries}m{0.17\textwidth}|m{\myhcolw}|},
    row{even}={gray!20}
  }
  \toprule
  Preparation: &
  The LA completes the end report using the template available in the \textit{Project templates} folder and submits it within three months after project completion.
  \\[1.5ex]
  Scientific review: &
  The Lead RSE reviews the scientific content, verifies that all project outputs have been registered in the appropriate systems, and requests revisions from the LA if needed.
  \\[1.5ex]
  Approval: &
  The responsible SH reviews and approves the end report before it is submitted to Finance for administrative processing.
  \\[1.5ex]
  Financial review: &
  Finance reviews the financial information, where applicable, and requests corrections if needed.
  \\[1.5ex]
  Formal closure: &
  Once all reviews have been completed, Finance archives the required documentation and issues the formal project closing letter to the LA, thereby formally concluding the project.
  \\[1.5ex]
  \bottomrule
 \end{booktabs}
\end{table}

For external projects, the reporting requirements of the funding organisation take precedence 
(see Appendix~\ref{app:external-projects}). Whenever possible, the report submitted to the funding organisation is 
also archived in the project portfolio.

The end report typically includes a public summary, project outcomes, challenges, opportunities, sustainability plans, 
project deliverables, workshops, and, where applicable, a financial overview. The completed report must be submitted 
within three months after the end of the project. 

After receiving the completed end report, the Lead RSE reviews the scientific content and ensures that all project 
deliverables and outputs have been registered in the appropriate systems. This includes updating the RSD project page, 
including the public summary and any missing project metadata, archiving the latest project documentation (e.g., the 
end report, Technology Plan, DMP and SMP, where applicable) in the project portfolio, and requesting revisions from the 
LA if necessary. Once the scientific review has been completed, the responsible SH approves the report and forwards it 
to Finance for review of the financial information. Communications may use the public summary and project outcomes for 
dissemination activities.

If the report is approved, Finance completes the administrative closure of the project, archives the required 
documentation in the project portfolio, and issues the formal project closing letter to the LA, thereby formally 
closing the project and ending the project lifecycle.